\section{Estimator and Likelihood Details}
\label{app:estimator}

This appendix separates two evaluated computations from the unexecuted common
estimator. Sections~\ref{sec:supp-state-recursion}--\ref{sec:supp-evaluated-event-bank}
specify the state recursion and load-support posterior that produced the reported
results. The complete statistical object is the candidate-conditioned joint
posterior and its normalized evidence; filtering, smoothing, or optimization is
the numerical route used to approximate it. Section~\ref{sec:supp-joint-design}
specifies that common state--event--integrity object, but its final integration
has not yet been evaluated.

\subsection{Bayesian Recursion and the State Hierarchy}
\label{sec:supp-state-recursion}

For a Markov state model, Bayesian filtering alternates prediction and correction \cite{sarkka2013bayesian}. More broadly, this construction treats the PMU task as a Bayesian inverse problem in which physical candidate models and latent states are compared through predictive evidence \cite{stuart2010inverse}:
\begin{align}
 p(\bm x_k\mid\bm y_{0:k-1})
 &=\int p(\bm x_k\mid\bm x_{k-1})
 p(\bm x_{k-1}\mid\bm y_{0:k-1})\,d\bm x_{k-1},
 \label{eq:supp-bayes-prediction}\\
 p(\bm x_k\mid\bm y_{0:k})
 &=\frac{p(\bm y_k\mid\bm x_k)
 p(\bm x_k\mid\bm y_{0:k-1})}
 {\int p(\bm y_k\mid\bm x_k)
 p(\bm x_k\mid\bm y_{0:k-1})\,d\bm x_k}.
 \label{eq:supp-bayes-correction}
\end{align}
The denominator in \eqref{eq:supp-bayes-correction} is the predictive evidence. Omitting it is harmless for a state mean under one fixed model, but not when physical candidates are compared.

B0 returns the equilibrium. B1 conditions the equilibrium prior on one PMU frame:
\begin{align}
 \bm K_{\mathrm{LMMSE}}
 &=\bm P_0\bm C^{\mathsf T}
 (\bm C\bm P_0\bm C^{\mathsf T}+\bm R)^{-1},\\
 \widehat{\delta\bm x}^{\mathrm{B1}}_k
 &=\bm K_{\mathrm{LMMSE}}(\bm y_k-\bm y_0).
 \label{eq:supp-b1-update}
\end{align}
B2 retains state and covariance between frames. For the linear Gaussian model, substituting Gaussian densities into \eqref{eq:supp-bayes-prediction}--\eqref{eq:supp-bayes-correction} yields
\begin{align}
 \bm m_k^-&=\bm A_d\bm m_{k-1}^+,
 &\bm P_k^-&=\bm A_d\bm P_{k-1}^+\bm A_d^{\mathsf T}+\bm Q,\label{eq:supp-kf-predict}\\
 \bm S_k&=\bm C\bm P_k^-\bm C^{\mathsf T}+\bm R,
 &\bm K_k&=\bm P_k^-\bm C^{\mathsf T}\bm S_k^{-1},\label{eq:supp-kf-gain}\\
 \bm m_k^+&=\bm m_k^-+\bm K_k(\bm y_k-\bm C\bm m_k^-),
 &\bm P_k^+&=(\bm I-\bm K_k\bm C)\bm P_k^-.
 \label{eq:supp-kf-correct}
\end{align}

\begin{theorem}[Linear-Gaussian reduction]
Equations~\eqref{eq:supp-kf-predict}--\eqref{eq:supp-kf-correct} are the exact posterior mean and covariance generated by the Bayesian recursion for the stated linear Gaussian model.
\end{theorem}

\paragraph{Proof.}
An affine transformation of the Gaussian posterior at $k-1$, followed by addition of independent Gaussian process noise, gives \eqref{eq:supp-kf-predict}. The pair $(\bm x_k,\bm y_k)$ is jointly Gaussian with cross covariance $\bm P_k^-\bm C^{\mathsf T}$ and measurement covariance $\bm S_k$. The standard conditional-Gaussian formula gives \eqref{eq:supp-kf-gain}--\eqref{eq:supp-kf-correct}. Equivalently, multiplying the predictive density by the likelihood and completing the square gives posterior precision $(\bm P_k^-)^{-1}+\bm C^{\mathsf T}\bm R^{-1}\bm C$; Woodbury's identity produces the same covariance and mean. $\hfill\square$

The implementation may use the Joseph covariance form
\begin{equation}
 \bm P_k^+=(\bm I-\bm K_k\bm C)\bm P_k^-
 (\bm I-\bm K_k\bm C)^{\mathsf T}+\bm K_k\bm R\bm K_k^{\mathsf T}
 \label{eq:supp-joseph-form}
\end{equation}
to preserve symmetry and positive semidefiniteness numerically.

B3 uses the Rauch--Tung--Striebel backward recursion \cite{rauch1965smoothing}:
\begin{align}
 \bm J_t&=\bm P_t^+\bm A_d^{\mathsf T}(\bm P_{t+1}^-)^{-1},\\
 \bm m_{t\mid k}&=\bm m_t^+
 +\bm J_t(\bm m_{t+1\mid k}-\bm m_{t+1}^-),\\
 \bm P_{t\mid k}&=\bm P_t^+
 +\bm J_t(\bm P_{t+1\mid k}-\bm P_{t+1}^-)
 \bm J_t^{\mathsf T}.
 \label{eq:supp-rts}
\end{align}
It estimates a past state using later measurements. The negative B3 result therefore diagnoses model mismatch; it does not imply that correct future information generally lacks value.

The nominal uncertainty correction rescales the rectangular voltage covariance as
\begin{equation}
 \bm\Sigma_{i,k}^{\mathrm{U2}}=\bm T\bm\Sigma_{i,k}\bm T^{\mathsf T},
 \qquad \bm T=\operatorname{diag}(\sqrt{\tau_{\Re}},\sqrt{\tau_{\Im}}),
 \label{eq:supp-temperature-map}
\end{equation}
with temperatures fitted on independent calibration trajectories. The observed-only adequacy statistic is a five-frame trailing mean of
\begin{equation}
 a_k=(\bm\nu_k-\bm\mu_{\nu})^{\mathsf T}
 \bm\Sigma_{\nu}^{-1}(\bm\nu_k-\bm\mu_{\nu}).
 \label{eq:supp-adequacy-statistic}
\end{equation}
Neither operation changes the physical state estimate.

\subsection{Constrained Physical Recentering}

The static recentering estimator treats active/reactive load changes and active generator changes as nuisance coordinates $\bm d$. It solves
\begin{equation}
 \begin{aligned}
 (\widehat{\bm V},\widehat{\bm d})
 =\arg\min_{\bm V,\bm d}\;&
 \|\bm e_y(\bm V)\|_2^2
 +\gamma^2\|\bm g_{\mathrm{AC}}(\bm V,\bm d)\|_2^2\\[-1mm]
 &+\lambda q_d^{-2}\|\bm d\|_2^2.
 \end{aligned}
 \label{eq:supp-recentering-objective}
\end{equation}
Development data fix $\gamma=10^3$, $q_d=1$, and $\lambda=0.01$. Because the AC constraint appears through a finite penalty rather than a normalized stochastic density, the implemented objective is a regularized physical estimator, not a fully normalized posterior.

For an exactly constrained local MAP, let $\bm r(\bm\chi)$ be the stacked residual, $\bm c(\bm\chi)=\bm0$ the equality constraints, and $\bm\Omega$ the residual precision. A Gauss--Newton step solves
\begin{equation}
 \begin{bmatrix}
 \bm H & \bm G_c^{\mathsf T}\\
 \bm G_c & \bm0
 \end{bmatrix}
 \begin{bmatrix}\Delta\bm\chi\\\bm\lambda\end{bmatrix}
 =
 \begin{bmatrix}-\nabla J\\-\bm c\end{bmatrix},
 \quad
 \bm H=\bm J_r^{\mathsf T}\bm\Omega\bm J_r+\bm H_{\mathrm{prior}}.
 \label{eq:supp-kkt-system}
\end{equation}

\begin{theorem}[Feasible local covariance]
If $\bm G_c$ has full row rank and $\bm Z$ spans $\ker\bm G_c$, then $\bm Z^{\mathsf T}\bm H\bm Z\succ\bm0$ makes the KKT matrix nonsingular. The local constrained Gaussian covariance is
\begin{equation}
 \bm P_{\chi}=\bm Z(\bm Z^{\mathsf T}\bm H\bm Z)^{-1}\bm Z^{\mathsf T}.
 \label{eq:supp-constrained-covariance}
\end{equation}
\end{theorem}

\paragraph{Proof.}
Any vector in the KKT nullspace satisfies $\bm G_c\bm v=\bm0$, hence $\bm v=\bm Z\bm q$. Premultiplying the first block equation by $\bm v^{\mathsf T}$ gives
$\bm q^{\mathsf T}\bm Z^{\mathsf T}\bm H\bm Z\bm q=0$, so $\bm q=\bm0$ and $\bm v=\bm0$. Full row rank then forces the multiplier to zero. Parameterizing feasible perturbations by $\delta\bm\chi=\bm Z\bm q$ gives Gaussian precision $\bm Z^{\mathsf T}\bm H\bm Z$ in $\bm q$; its linear pushforward is \eqref{eq:supp-constrained-covariance}. $\hfill\square$

This covariance is local and unimodal. The reported joint NEES-like failure shows that it is not calibrated for the full shifted-plant experiment.

\subsection{Evaluated Physical Load-Support Model}
\label{sec:supp-evaluated-event-bank}

Let $\bm r$ be the centered $T$-frame PMU window. For support $S$ with amplitudes $\bm a_S$, the evaluated mean is
\begin{equation}
 \bm\mu_S(\bm a_S)=
 \sum_{i\in S}(a_i\bm D_i+a_i^2\bm Q_i)
 +\sum_{\substack{i,j\in S\\i<j}}a_ia_j\bm Q_{ij}.
 \label{eq:supp-evaluated-event-mean}
\end{equation}
The W2 likelihood is
\begin{align}
 q_S(\bm a_S)
 &=\|\bm r-\bm\mu_S(\bm a_S)\|_{\bm\Sigma_T^{-1}}^2,\\
 p(\bm r\mid S,\bm a_S)
 &=\frac{\exp[-q_S(\bm a_S)/2]}
 {(2\pi)^{n_yT/2}|\bm\Sigma_T|^{1/2}},\\
 \bm\Sigma_T&=\bm R_T(\rho)\otimes\bm\Omega.
 \label{eq:supp-evaluated-likelihood}
\end{align}
The covariance is fitted only on normal calibration records and remains fixed across supports. The selected priors are
\begin{align}
 p(K=0,1,2)&=(0.20,0.50,0.30),\\
 p(S\mid K)&=\binom{16}{K}^{-1},\\
 a_i\mid i\in S&\sim\mathcal N(0,\sigma_a^2),
 &\sigma_a&=0.05.
 \label{eq:supp-evaluated-priors}
\end{align}

For reference, the linear single-source model has a closed form.

\begin{proposition}[Conjugate amplitude posterior and Bayes factor]
Let $\bm r=a\bm d+\bm\varepsilon$, $\bm\varepsilon\sim\mathcal N(\bm0,\bm\Sigma)$, and $a\sim\mathcal N(0,\sigma_a^2)$. Define
$I=\bm d^{\mathsf T}\bm\Sigma^{-1}\bm d$ and
$s=\bm d^{\mathsf T}\bm\Sigma^{-1}\bm r$. Then
\begin{align}
 v&=(I+\sigma_a^{-2})^{-1},
 &m&=vs,\\
 a\mid\bm r,S&\sim\mathcal N(m,v),\\
 \log\frac{p(\bm r\mid S)}{p(\bm r\mid H_0)}
 &=\tfrac12\frac{s^2}{I+\sigma_a^{-2}}
 -\tfrac12\log(1+\sigma_a^2I).
 \label{eq:supp-linear-bayes-factor}
\end{align}
\label{prop:supp-conjugate-bayes}
\end{proposition}

\paragraph{Proof.}
Terms in the joint log density that depend on $a$ are
$-\tfrac12[(I+\sigma_a^{-2})a^2-2sa]$.
Completing the square gives mean $m$ and variance $v$. Integrating the resulting Gaussian contributes
$\sqrt{v/\sigma_a^2}\exp(vs^2/2)$ relative to $H_0$; taking its logarithm gives \eqref{eq:supp-linear-bayes-factor}. The logarithmic term is the amplitude-volume or Occam penalty. $\hfill\square$

\begin{corollary}[One-sided amplitude evidence]
Under Proposition~\ref{prop:supp-conjugate-bayes}, replace the signed Gaussian
slab by the proper half-normal prior
$p(a)=2\mathcal N(a;0,\sigma_a^2)\mathbb I(a\ge0)$. If
$BF_{\pm}$ denotes the signed Bayes factor in
\eqref{eq:supp-linear-bayes-factor}, then
\begin{equation}
 BF_{+}=2\Phi(z)BF_{\pm},
 \qquad
 z=\frac{\sigma_a s}{\sqrt{1+\sigma_a^2I}}.
 \label{eq:supp-one-sided-bayes-factor}
\end{equation}
Conditional on the active model, the severity posterior is
$\mathcal N(m,v)$ truncated to $[0,\infty)$.
\end{corollary}

\paragraph{Proof.}
The fixed-amplitude likelihood ratio is
$\exp(as-Ia^2/2)$. Multiplication by the half-normal density and completion of
the square give twice the signed Gaussian integral restricted to $a\ge0$.
The unrestricted posterior has standardized probability
$\Phi(m/\sqrt v)=\Phi(z)$ on that half-line, which yields
\eqref{eq:supp-one-sided-bayes-factor}. $\hfill\square$

The factor $2\Phi(z)$ preserves direction: a large score with the prohibited
sign is evidence against the one-sided event even though the signed factor
depends on $s^2$. Proper slab scales are essential because Bayes factors across
different support dimensions are undefined up to an arbitrary constant under
an improper amplitude prior. Under the one-sided null $a=0$, the likelihood
ratio statistic satisfies
\begin{equation}
 2\Delta\ell=[Z]_+^2,\qquad Z\sim\mathcal N(0,1),
 \label{eq:supp-boundary-lrt}
\end{equation}
and therefore has the boundary mixture
$\tfrac12\delta_0+\tfrac12\chi_1^2$, rather than the regular $\chi_1^2$ law.

The quadratic multi-source mean is nonconjugate. Its evidence is
\begin{equation}
 Z_S=p(\bm r\mid S)=
 \int p(\bm r\mid S,\bm a_S)
 \prod_{i\in S}p(a_i)\,d\bm a_S.
 \label{eq:supp-support-evidence}
\end{equation}
With $Q$-node Gauss--Hermite abscissas $x_q$ and weights $w_q$, the Gaussian-prior transformation $a_i=\sqrt2\sigma_ax_{q_i}$ gives
\begin{equation}
 Z_S\approx\pi^{-|S|/2}
 \sum_{q_1=1}^{Q}\cdots\sum_{q_{|S|}=1}^{Q}
 \left(\prod_{j=1}^{|S|}w_{q_j}\right)
 p(\bm r\mid S,\sqrt2\sigma_a\bm x_{\bm q}).
 \label{eq:supp-gauss-hermite-evidence}
\end{equation}
The implementation uses $Q=31$ and evaluates \eqref{eq:supp-gauss-hermite-evidence} in the log domain using log-sum-exp. Independent adaptive Gauss--Kronrod integration checks representative two-dimensional cases. These checks support numerical integration under the frozen model; they do not validate the physical model itself.

\subsection{Posterior Outputs and the Cardinality--Localization Distinction}

The normalized support posterior is
\begin{equation}
 p(S\mid\bm r)=
 \frac{Z_Sp(S)}{\sum_{R\in\mathcal H}Z_Rp(R)}.
 \label{eq:supp-support-posterior}
\end{equation}
It induces cardinality, source-inclusion, and amplitude distributions:
\begin{align}
 p(K=k\mid\bm r)&=\sum_{|S|=k}p(S\mid\bm r),\label{eq:supp-cardinality-posterior}\\
 p(i\in S\mid\bm r)&=\sum_{S:\,i\in S}p(S\mid\bm r),\label{eq:supp-inclusion-posterior}\\
 p(a_i\mid\bm r)&=p(i\notin S\mid\bm r)\delta_0
 +\sum_{S:\,i\in S}p(S\mid\bm r)p(a_i\mid S,\bm r).
 \label{eq:supp-spike-slab-severity}
\end{align}
Equation~\eqref{eq:supp-spike-slab-severity} is a model-averaged spike-and-slab output: source absence contributes a point mass at zero.

The experimentally observed gap between cardinality and support has an exact algebraic decomposition. Let $N_k=\binom{16}{k}$, $\ell_S=\log Z_S$, and $\ell_k^{\star}=\max_{|S|=k}\ell_S$. Define the effective multiplicity
\begin{equation}
 m_{\mathrm{eff},k}=\sum_{|S|=k}\exp(\ell_S-\ell_k^{\star}).
 \label{eq:supp-effective-multiplicity}
\end{equation}

\begin{proposition}[Exact effective-multiplicity identity]
Under a uniform prior over supports of cardinality $k$,
\begin{align}
 \log p(\bm r\mid K=k)
 &=\ell_k^{\star}-\log N_k+\log m_{\mathrm{eff},k},
 \label{eq:supp-cardinality-evidence-decomposition}\\
 \max_{|S|=k}p(S\mid K=k,\bm r)
 &=m_{\mathrm{eff},k}^{-1}.
 \label{eq:supp-max-support-multiplicity}
\end{align}
\end{proposition}

\paragraph{Proof.}
Factor $\exp(\ell_k^{\star})$ from
$N_k^{-1}\sum_{|S|=k}\exp(\ell_S)$ and take logarithms to obtain \eqref{eq:supp-cardinality-evidence-decomposition}. Conditional normalization within cardinality $k$ divides the largest evidence by the same sum, giving \eqref{eq:supp-max-support-multiplicity}. $\hfill\square$

Many similar pair supports can therefore increase evidence for $K=2$ while preventing any one pair from dominating. Entropy separates the same two questions exactly because $K$ is a deterministic function of $S$:
\begin{equation}
 H(S\mid\bm r)=H(K\mid\bm r)
 +\sum_kp(K=k\mid\bm r)H(S\mid K=k,\bm r).
 \label{eq:supp-support-entropy-decomposition}
\end{equation}
The first term measures cardinality uncertainty; the second measures localization ambiguity within each cardinality.

For set-valued output, sort supports by posterior mass and define the smallest credible set
\begin{equation}
 \mathcal C_q(\bm r)=\arg\min_{\mathcal C}|\mathcal C|
 \quad\text{subject to}\quad
 \sum_{S\in\mathcal C}p(S\mid\bm r)\ge q.
 \label{eq:supp-credible-support-set}
\end{equation}
Candidate-set size and coverage are more informative than an arbitrary Top-$m$ when the posterior is multimodal. The present paper has not yet calibrated an abstention threshold or credible-set coverage prospectively.

\subsection{Step-by-Step Evaluated Bayesian Algorithm}

Table~\ref{tab:supp-evaluated-algorithm} records the exact order of the load-support computation. Steps marked ``evaluated'' produced the 137-hypothesis results; optional outputs are mathematically defined but are not promoted to validated claims.

\begin{table*}[t]
\caption{Evaluated 137-hypothesis Bayesian load-support algorithm.}
\label{tab:supp-evaluated-algorithm}
\centering
\scriptsize
\setlength{\tabcolsep}{4pt}
\begin{tabularx}{\textwidth}{@{}c l X l@{}}
\toprule
Step & Operation & Computation and frozen choice & Status \\
\midrule
0 & Declare contract & Fix 16 candidate load buses, $K\le2$, known onset, 30-frame window, eight PMUs, nominal topology, and no corrupted channels & Evaluated \\
1 & Validate plant & Check equilibrium, terminal-current orientation, sampling, nonlinear solve acceptance, and first-order response before building the bank & Evaluated \\
2 & Build physical dictionary & Generate centered $\bm D_i$, self-curvature $\bm Q_i$, and cross-interaction $\bm Q_{ij}$ from simulator-consistent interventions & Evaluated \\
3 & Fit noise on CAL & Estimate channel covariance and AR(1) dependence from normal calibration records; freeze W2 before event TEST & Evaluated \\
4 & Freeze priors and integrator & Set $p(K)$, uniform support priors, Gaussian amplitude slab, GH31 nodes/weights, and independent GK checks & Evaluated \\
5 & Form the residual & Align the known onset, subtract the frozen normal mean, retain the fixed PMU/channel order, and apply the declared temporal covariance & Evaluated \\
6 & Enumerate supports & Construct $H_0$, 16 single-source supports, and 120 two-source supports; no learned source label prunes the bank & Evaluated \\
7 & Integrate each support & Evaluate \eqref{eq:supp-support-evidence} with \eqref{eq:supp-gauss-hermite-evidence}; retain quadrature-weighted amplitude moments & Evaluated \\
8 & Normalize once & Add the support prior in log space, subtract the global log-sum-exp, and obtain \eqref{eq:supp-support-posterior} & Evaluated \\
9 & Read decisions in order & Report event presence, $p(K\mid\bm r)$, conditional support ranking, source inclusion, and conditional amplitude intervals separately & Evaluated \\
10 & Quantify ambiguity & Compute entropy, effective multiplicity, or $\mathcal C_q$ without relabeling them as validated abstention guarantees & Defined; partial evidence \\
11 & Record provenance & Store the plant identifier, support, amplitude, noise seed, likelihood version, and frozen artifact hash for every independent record & Evaluated \\
\bottomrule
\end{tabularx}
\end{table*}

The direct quadrature cost per record is
\begin{equation}
 1+GQ+\binom{G}{2}Q^2
 \label{eq:supp-gh-complexity}
\end{equation}
likelihood-node evaluations for $K\le2$. With $G=16$ and $Q=31$, this is 115,817 nodes before vectorization or reuse. The count explains why the support bank, integration order, and horizon must be frozen together when runtime is compared.

\subsection{Candidate-Conditioned Joint State--Event Design}
\label{sec:supp-joint-design}

The common-contract extension would infer state trajectory $\bm X$, algebraic variables $\bm Z$, event support $S$, amplitudes $\bm a_S$, operating nuisance $\bm\eta$, and integrity corruption $\bm b$. Its posterior factorization is
\begin{align}
 &p(\bm X,\bm Z,\bm\eta,\bm b,S,\bm a_S\mid\bm Y)
 \propto p(S)p(\bm a_S\mid S)p(\bm\chi_0)\nonumber\\
 &\quad\times\prod_{t=0}^{T-1}
 p(\bm\chi_{t+1}\mid\bm\chi_t,S,\bm a_S)
 \prod_{t=0}^{T}p(\bm y_t\mid\bm\chi_t,\bm b_t,S,\bm a_S),
 \label{eq:supp-full-joint-posterior}
\end{align}
subject to the candidate DAE constraints. Here $\bm\chi_t$ denotes a consistent state/nuisance chart. A constrained negative twice-log density, excluding $p(S)$, is
\begin{equation}
 \begin{aligned}
 \mathcal J_S^{\circ}={}&
 \|\bm\chi_0-\overline{\bm\chi}_0\|_{\bm P_0^{-1}}^2
 +\sum_{t=0}^{T}\|\bm r_{S,t}^{\mathrm{pmu}}\|_{\bm R_t^{-1}}^2\\
 &+\sum_{t=0}^{T-1}
 \|\bm r_{S,t}^{\mathrm{dyn}}\|_{\bm Q_t^{-1}}^2
 +2\lambda_b\sum_{t=0}^{T}\|\bm b_t\|_1.
 \end{aligned}
 \label{eq:supp-joint-objective}
\end{equation}
Candidate evidence must be evaluated in a common feasible measure. The boundary
correction is exposed by an explicit concentration parameter rather than by
informally assuming a concentrated posterior.

\begin{theorem}[Local cone--Laplace evidence]
Let $L_{S,\epsilon}$ be the log likelihood plus the log of a proper prior in
identifiable local coordinates of dimension $d_S$. Suppose the rescaled
feasible domains converge locally to a closed polyhedral cone $\bm K_S$ and
\begin{equation}
 L_{S,\epsilon}(\bm\xi_0+\epsilon\bm u)
 -L_{S,\epsilon}(\bm\xi_0)
 =\bm h_S^{\mathsf T}\bm u
 -\tfrac12\bm u^{\mathsf T}\bm H_S\bm u
 +r_\epsilon(\bm u),
 \label{eq:supp-cone-laplace-expansion}
\end{equation}
where $\bm H_S\succ\bm0$ and $r_\epsilon\to0$ uniformly on compact sets.
Assume an integrable dominating envelope and negligible evidence from distant
modes. Then
\begin{align}
 Z_{S,\epsilon}
 ={}&e^{L_{S,\epsilon}(\bm\xi_0)}\epsilon^{d_S}
 (2\pi)^{d_S/2}|\bm H_S|^{-1/2}
 \nonumber\\
 &\times\exp\!\left(
 \tfrac12\bm h_S^{\mathsf T}\bm H_S^{-1}\bm h_S
 \right)
 \Pr(\bm U_S\in\bm K_S)[1+o(1)],
 \label{eq:supp-laplace-evidence}\\
 \bm U_S\sim{}&
 \mathcal N(\bm H_S^{-1}\bm h_S,\bm H_S^{-1}).
 \nonumber
\end{align}
\label{thm:supp-cone-laplace}
\end{theorem}

\paragraph{Proof.}
Set $\bm\xi=\bm\xi_0+\epsilon\bm u$, which contributes the Jacobian
$\epsilon^{d_S}$. The compact-uniform expansion gives pointwise convergence of
the rescaled integrand to
$\exp(\bm h_S^{\mathsf T}\bm u-\bm u^{\mathsf T}\bm H_S\bm u/2)$.
The assumed envelope permits dominated convergence, while the distant-mode
condition localizes the original integral. Completing the square in the
limiting quadratic integral gives the Gaussian normalizer and the cone
probability in \eqref{eq:supp-laplace-evidence}. $\hfill\square$

For an interior model $\bm K_S=\mathbb R^{d_S}$ and the probability factor is
one. For one or more boundary amplitudes it is a Gaussian orthant probability,
not generally $2^{-q}$. Half-normal prior normalizers supply their own factors
in addition to this cone probability. Chart Jacobians, prior normalizers, and
candidate-dependent covariance determinants belong in $L_{S,\epsilon}$ or the
measure term. Comparing only optimized residuals drops both volume and boundary
terms and does not produce a posterior.

For a small-noise Gaussian Hybrid-DAE trajectory map
$\bm\mu_S(\bm\xi_0+\epsilon\bm u)
=\bm\mu_0+\epsilon\bm J_S\bm u+O(\epsilon^2)$,
\begin{equation}
 \bm h_S=\bm J_S^{\mathsf T}\bm\Sigma^{-1}\bm z,\qquad
 \bm H_S=\bm J_S^{\mathsf T}\bm\Sigma^{-1}\bm J_S,
 \label{eq:supp-dae-cone-laplace-hessian}
\end{equation}
so positive curvature is equivalent to first-order identifiability in the
reduced coordinates. This is a theoretical local specialization, not a
validation of the unevaluated joint estimator. If $\bm H_S$ is singular, merely
replacing its inverse by a pseudoinverse is invalid: redundant coordinates must
first be quotiented out, while genuine higher-order identification changes the
evidence scale. For example,
\begin{equation}
 \int_0^\infty e^{-na^4/4}\,da
 =\frac{\Gamma(1/4)}{4^{3/4}}n^{-1/4},
 \label{eq:supp-singular-quartic-evidence}
\end{equation}
not the $n^{-1/2}$ rate of regular quadratic Laplace theory.

\begin{proposition}[Exact reduction along redundant fibers]
Suppose a parameter chart $\bm\theta\leftrightarrow(\bm\xi,\bm\zeta)$ is
one-to-one on the prior support and the observable map satisfies
$\bm\mu(\bm\theta)=\overline{\bm\mu}(\bm\xi)$ there. If
$\widetilde\pi(\bm\xi,\bm\zeta)$ is the transformed proper prior, then the
evidence equals the reduced integral with induced prior
\begin{equation}
 \overline\pi(\bm\xi)=
 \int\widetilde\pi(\bm\xi,\bm\zeta)\,d\bm\zeta.
 \label{eq:supp-induced-fiber-prior}
\end{equation}
Thus an exactly redundant coordinate contributes no vanishing small-noise
power; its effect is the pushforward prior measure, not a pseudodeterminant.
\label{prop:supp-redundant-fiber}
\end{proposition}

\paragraph{Proof.}
The likelihood is independent of $\bm\zeta$. Fubini's theorem integrates the
proper transformed prior along each fiber and gives
$Z=\int p(\bm y\mid\bm\xi)\overline\pi(\bm\xi)d\bm\xi$ exactly.
$\hfill\square$

For a merely local chart, the same reduction is asymptotic only after a
localization condition excludes evidence from distant branches.

\begin{theorem}[Mixed regular and higher-order evidence]
In whitened coordinates, suppose the first observable expansion after
redundancy removal is
\begin{equation}
 \bm\mu(\bm\xi,a)=\bm\mu_0+\bm A\bm\xi
 +\frac{\bm c_k}{k!}a^k
 +o(\|\bm\xi\|+|a|^k),
 \label{eq:supp-mixed-singular-map}
\end{equation}
where $\bm A$ has rank $r$, the proper prior density is continuous and positive
at the origin, and
\begin{equation}
 \bm c_{k,\perp}=(\bm I-\bm P_{\bm A})\bm c_k\ne\bm0.
 \label{eq:supp-normal-higher-order-coefficient}
\end{equation}
Assume a dominating envelope and negligible evidence away from this branch.
For the Gaussian-normalizer-free evidence kernel with noise scale $\epsilon$,
\begin{align}
 M_\epsilon\sim{}&C\epsilon^{r+1/k},
 \label{eq:supp-mixed-singular-evidence-rate}\\
 C={}&\pi(\bm0,0)(2\pi)^{r/2}
 |\bm A^{\mathsf T}\bm A|^{-1/2}\\
 &\times
 \frac{\Gamma(1/(2k))}{k}
 \left[\frac{2(k!)^2}{\|\bm c_{k,\perp}\|_2^2}
 \right]^{1/(2k)}.
 \label{eq:supp-mixed-singular-evidence-constant}
\end{align}
\label{thm:supp-mixed-singular-evidence}
\end{theorem}

\paragraph{Proof.}
Use $\bm\xi=\epsilon\bm u$ and
$a=\epsilon^{1/k}t$. The Jacobian is
$\epsilon^{r+1/k}$ and the limiting kernel is
$\exp[-\|\bm A\bm u+\bm c_kt^k/k!\|^2/2]$.
The component of $\bm c_k$ in $\operatorname{range}\bm A$ disappears after a
translation of $\bm u$. Gaussian integration over the remaining regular
coordinates yields $(2\pi)^{r/2}|\bm A^{\mathsf T}\bm A|^{-1/2}$; integration
of the normal component over $t\in\mathbb R$ gives the gamma-function factor.
Dominated convergence supplies the prior density at the origin and the stated
limit. $\hfill\square$

For one additional amplitude, the theorem separates three cardinality cases.
A regular direction supplies an $O(\epsilon)$ Occam factor, a direction first
visible at order $k$ supplies $O(\epsilon^{1/k})$, and an exactly redundant
extension supplies an $O(1)$ factor determined by its induced prior. For
$\epsilon\asymp n^{-1/2}$ these become $n^{-1/2}$,
$n^{-1/(2k)}$, and $n^0$. These are local evidence rates; the RAWSIM39 gate
ledger does not contain the higher-order derivatives required to assign them
to its rejected parameters.

\begin{table}[t]
\caption{Deterministic checks of boundary-aware evidence identities. These toy systems are not evaluated estimator results.}
\label{tab:supp-boundary-bayes-checks}
\centering
\scriptsize
\begin{tabularx}{\columnwidth}{@{}p{0.40\columnwidth}X@{}}
\toprule
Check & Recomputed outcome \\
\midrule
Signed/one-sided Bayes factors & max quadrature error 2.7e-12 (pass). \\
Broad nuisance limit & relative information error 8.9e-10 (pass). \\
Boundary likelihood ratio & zero mass 0.4972; positive mean 0.9987 (pass). \\
Nonlinear cone--Laplace & relative-error slope 0.9497 (pass). \\
Singular quartic evidence & scaling slope -0.2500 (pass). \\
\bottomrule
\end{tabularx}
\end{table}


For a sliding window, eliminating old Gaussian variables creates an arrival prior. Write its canonical log kernel as
$-\tfrac12\bm z^{\mathsf T}\bm\Lambda\bm z+\bm h^{\mathsf T}\bm z+c$.
If the information matrix and vector are partitioned into variables $\bm a$ to eliminate and retained variables $\bm b$, exact integration yields
\begin{align}
 \bm\Lambda_{\mathrm{new}}
 &=\bm\Lambda_{bb}-\bm\Lambda_{ba}\bm\Lambda_{aa}^{-1}\bm\Lambda_{ab},\label{eq:supp-schur-precision}\\
 \bm h_{\mathrm{new}}
 &=\bm h_b-\bm\Lambda_{ba}\bm\Lambda_{aa}^{-1}\bm h_a,\label{eq:supp-schur-vector}\\
 c_{\mathrm{new}}
 &=c+\tfrac12\bm h_a^{\mathsf T}\bm\Lambda_{aa}^{-1}\bm h_a
 +\tfrac{n_a}{2}\log(2\pi)-\tfrac12\log\det\bm\Lambda_{aa}.
 \label{eq:supp-schur-normalizer}
\end{align}
The first two lines preserve the next MAP problem; the last line is also required to preserve evidence between candidates.

Table~\ref{tab:supp-joint-algorithm} gives the intended sequence. It is a specification for the next experiment, not a claim about the results reported in the main paper.

\begin{table*}[t]
\caption{Candidate-conditioned joint estimator design and current implementation boundary.}
\label{tab:supp-joint-algorithm}
\centering
\scriptsize
\setlength{\tabcolsep}{4pt}
\begin{tabularx}{\textwidth}{@{}c l X l@{}}
\toprule
Step & Operation & Required calculation & Present status \\
\midrule
1 & Pre-event update & Use only pre-event measurements to update state and operating nuisance; retain cross covariance and normalized arrival evidence & State filter evaluated; joint nuisance pending \\
2 & Integrity mask & Mark missing channels exactly and introduce sparse corruption variables for available but suspect channels & Ontology defined; common inference pending \\
3 & Candidate generation & Enumerate physical family, support, onset, and severity domain without using source identity as a learned label & Load family evaluated; other families pending \\
4 & Candidate prediction & Apply the candidate DAE/event map and the current PMU membership to generate a physical trajectory & Load bank evaluated \\
5 & Candidate-conditioned solve & Minimize \eqref{eq:supp-joint-objective} on the feasible DAE manifold using the KKT system and continuation from the pre-event state & Not executed jointly \\
6 & Learned discrepancy & If used, predict only candidate-residual discrepancy and covariance; do not replace the physical hypothesis or consume target source labels & Not implemented \\
7 & Evidence normalization & Retain log determinants, prior densities, feasible-coordinate measures, and the Schur normalizer before comparing candidates & GH31 load evidence evaluated; joint Laplace evidence pending \\
8 & Posterior outputs & Return state-functional mixture, cardinality, support, integrity state, severity, and candidate-set uncertainty & Separate state/load outputs only \\
9 & Sequential decision & Continue observing when expected information gain justifies delay; use distinct thresholds for detection, cardinality, and exact source & Not evaluated \\
10 & Reconfigure and verify & After accepting a physical change, rebuild the operating center and check subsequent innovations before declaring a persistent model update & Static recentering only \\
\bottomrule
\end{tabularx}
\end{table*}

For a posterior state $\bm\pi_t(S)=p(S\mid\bm Y_{1:t})$, an idealized unit-cost support decision with per-frame delay cost $c$ satisfies the Bellman equation
\begin{equation}
 V(\bm\pi)=\min\left\{
 1-\max_S\pi(S),\;
 c+\mathbb E[V(\bm\pi_{t+1})\mid\bm\pi_t=\bm\pi]
 \right\}.
 \label{eq:supp-bayes-stopping}
\end{equation}
The one-step expected information gain is
\begin{equation}
 I(S;\bm y_{t+1}\mid\bm Y_{1:t})
 =H(S\mid\bm Y_{1:t})
 -\mathbb E[H(S\mid\bm Y_{1:t+1})].
 \label{eq:supp-next-frame-information}
\end{equation}
Equations~\eqref{eq:supp-bayes-stopping}--\eqref{eq:supp-next-frame-information} motivate adaptive horizon selection. They are not an executed stopping rule for the current fixed-onset, fixed-window bank.

Finally, the hidden-voltage posterior must average across support uncertainty rather than report only the winning support. With $w_S=p(S\mid\bm Y)$, conditional target mean $\bm m_S$, and covariance $\bm P_S$,
\begin{align}
 \overline{\bm m}&=\sum_Sw_S\bm m_S,\label{eq:supp-mixture-mean}\\
 \operatorname{Cov}(\bm\psi\mid\bm Y)
 &=\sum_Sw_S\bm P_S
 +\sum_Sw_S(\bm m_S-\overline{\bm m})
 (\bm m_S-\overline{\bm m})^{\mathsf T}.
 \label{eq:supp-total-mixture-covariance}
\end{align}
The second term is between-support ambiguity. Dropping it would make the virtual-PMU field appear overconfident precisely when source hypotheses disagree.

\subsection{Estimator Claim Boundary}

The executed state recursion is Bayesian under its local linear Gaussian model, and the load-source bank computes normalized Bayesian evidence under the frozen W2/L2 model. The physical recentering is a regularized constrained estimator with separately assessed uncertainty. The full candidate-conditioned state--event--integrity posterior, learned discrepancy, unknown-onset inference, changing-PMU operation, and sequential stopping rule remain specified but unexecuted. No result in the main paper is attributed to those pending components.
